\documentclass[11pt]{article}
\usepackage{fontspec}
\usepackage{amsmath,amssymb,amsthm,mathtools}
\usepackage[margin=1.1in]{geometry}
\usepackage{microtype}
\usepackage{parskip}
\usepackage{lmodern}
\usepackage[colorlinks=true,linkcolor=blue!50!black,urlcolor=blue!50!black]{hyperref}
\providecommand{\tightlist}{}
\title{Reflection positivity and the obstruction to finite hidden-Markov realization}
\author{Felix Robles Elvira (ORCID: 0009-0009-2017-4394; independent researcher)}
\date{\small preprint, not peer reviewed, version 2026-06-11.}
\begin{document}
\maketitle

\section*{Abstract}
A stationary binary process is reflection-positive (RP) if the Gram matrices built from word probabilities by reflection through bonds ($G[u,v] = p(\bar u \circ v)$, $\bar u$ the reversal) \emph{and} through sites ($p(\bar u\,x\,v)$) are positive semidefinite — the discrete Osterwalder–Schrader structure, with the two reflections playing genuinely different roles.  We prove that RP processes are exactly those with a representation $p(x_1\cdots x_n) = \langle\omega, A_{x_1}\cdots A_{x_n}\omega\rangle$ with self-adjoint PSD letters summing to a contraction fixing $\omega$ — the POVM-type "quantum" realizations (site reflection is precisely what makes the letters positive) — and that within this class \emph{reflection positivity is exhausted}: every reflected Gram is automatically PSD, so the entire content of the open realization problem resides in word positivity of the non-commutative form.  Our main structural results confine the known obstruction to classical (entrywise-nonnegative, hidden-Markov) realizability: (i) for every block $W$ whose cyclic word is achiral, the diagonal $p(W^n)$ is a Stieltjes moment sequence at palindromic phase \emph{where such a phase exists} (odd length always; even length iff an edge–edge axis) and has real nonnegative poles at every phase, hence is positively realizable by the dominant-pole criterion; consequently any irrational-rotation ("clock") obstruction must live on \emph{chiral} blocks of length $\ge 6$; (ii) for the minimal dangerous configuration — a coupled spectral triple $\{\rho, \rho e^{\pm i\theta}\}$ on a chiral block — all constraints from the block–reversed-block word algebra reduce, exactly and $\theta$-freely, to torus-positivity conditions on a six-parameter reduced datum, and the \emph{normal point} of that datum satisfies every such condition at every coupling ratio $r < 1$ (the supremum $r = 1$ is approached but not attained), so no exclusion argument can be carried by the block algebra alone; (iii) we reduce the remaining conjecture, \emph{for rank-three blocks} (the higher-rank case needs a named strengthening stated in Section 6), to two sharp finite-dimensional statements — "normal words in two PSD letters have real spectrum" and "all-depth torus feasibility is isolated at the normal point" — and supply quantitative evidence including a two-walls separation: letter-realizable near-triples live at Gram distance $\ge 0.50$ from normality while deep-torus-feasible data live at $\le 0.15$.  An explicit three-dimensional reversible process whose diagonal oscillates on its decay circle (a "clock") shows the finite-rank realization question has genuine content; the clock provably fails reflection positivity.

\section*{1. Introduction}
\textbf{The realization problem.}  When is a stationary process a function of a finite-state Markov chain?  The question is classical (Dharmadhikari [1]; Heller's cone-theoretic characterization [2]) and hard: finite Hankel rank — the existence of a finite-dimensional \emph{quasi-realization} $p(w) = \pi M_{x_1}\cdots M_{x_n}\mathbf 1$ with sign-indefinite data — does not suffice.  The obstruction is spectral: Anderson-type positive-realization theory [3, 4] shows a diagonal sequence with a unique positive dominant pole is positively realizable, while an irrational rotation on the decay circle ($h(n) = \rho^n(a + b\cos(n\theta + \varphi))$, $\theta/\pi$ irrational) admits no finite positive realization at any dimension. We call a process exhibiting this on some block diagonal a \textbf{clock}.  Clocks exist: Section 3 presents a rank-three, \emph{reversible}, fully valid example, certified analytically.

\textbf{Reflection positivity.}  Bond-RP — positivity of the reversal-pairing Grams — is the probabilistic shadow of Osterwalder–Schrader positivity and of the transfer-matrix self-adjointness of equilibrium statistical mechanics [5].  The conjecture animating this paper, suggested by positive experience with reflection-positive systems:

\begin{quote}
\textbf{Conjecture (RP-realization).}  Reflection positivity (the bond-plus-site conjunction of Definition 2.0) plus finite Hankel rank implies a finite positive (hidden-Markov) realization.
\end{quote}
Equivalently: \emph{no clock hides inside a reflection-positive process.} The clock of Section 3 fails RP, consistent with the conjecture; the question is whether that is a theorem.

\textbf{Results.}  Section 2 proves the representation theorem (bond-RP $\iff$ POVM form) and an exhaustion principle: within the form, all reflected Grams are \emph{automatically} PSD — receipts include an engineered process with negative word probabilities whose reflected Grams are PSD to machine precision — so reflection contributes nothing beyond the form, and word positivity carries the entire problem.  Section 4 proves the confinement theorems via the companion paper's achiral-necklace theorem: achiral-block diagonals have exact Stieltjes structure at palindromic phase, real nonnegative poles at every phase, and are positively realizable; clocks are confined to chiral blocks of length $\ge 6$.  Section 5 analyzes the minimal dangerous configuration through an exact \emph{peripheral reduction} and proves the \emph{normal escape}: the block–reversed-block constraint tower is satisfiable at full coupling by normal-block data, closing every exclusion route that stays inside that algebra.  Section 6 states the resulting two-conjecture reduction and the evidence, including walls measured by three independent searches and a two-coordinate "pinch" showing the letter-realizable and torus-feasible regions are disjoint at every budget reached. Section 7 discusses the assembly problem and literature.

\section*{2. The representation theorem and the exhaustion of reflection}
\textbf{Definition 2.0 (reflection positivity, airtight).}  A stationary process $p$ on $\{0,1\}^{\mathbb Z}$ is \emph{reversible} if $p(\bar w) = p(w)$ for every word ($\bar w$ the reversal).  It is \textbf{bond-RP} if for every finite family of words $\{u_i\}$ the matrix $G[u_i, u_j] = p(\bar u_i \circ u_j)$ is PSD — where PSD is taken in the convention that \emph{includes symmetry of the matrix}, so reversibility ($p(\bar w) = p(w)$: take one word empty) is built into bond-RP by that convention rather than derived; we state reversibility separately in the definition for transparency.  (An earlier draft claimed reversibility follows "by polarization"; under the quadratic-form-only reading of PSD, polarization constrains only the symmetric part, so that claim is withdrawn and the convention is now explicit.)  It is \textbf{site-RP} if for every finite family and every letter $x$ the matrix $p(\bar u_i \, x \, u_j)$ is PSD.  We call $p$ \textbf{RP} if it is both bond- and site-RP — the conjunction is the discrete Osterwalder–Schrader structure (reflections through bonds \emph{and} through sites), and the two are genuinely different conditions (Section 2's theorem and remark make the division of labor precise).

\textbf{Theorem 2.1 (RP form).}  A stationary process $p$ is RP (Definition 2.0) iff there exist a Hilbert space, a \emph{unit} vector $\omega$, and self-adjoint operators $A_0, A_1 \succeq 0$ with $T = A_0 + A_1 \preceq \mathbf 1$, $T\omega = \omega$, such that $p(x_1\cdots x_n) = \langle\omega, A_{x_1}\cdots A_{x_n}\omega\rangle$. If the Hankel rank is finite the space may be taken finite-dimensional.

\emph{Proof.}  ($\Leftarrow$)  With $B_w = A_{x_1}\cdots A_{x_n}$, self-adjointness gives $B_{\bar w} = B_w^{\dagger}$, so the bond Gram is $\langle B_{u_i}\omega, B_{u_j}\omega\rangle$ and the site matrix is $\langle B_{u_i}\omega, A_x B_{u_j}\omega\rangle$ — PSD since $A_x \succeq 0$.  ($\Rightarrow$)  GNS/OS construction at a bond, run \emph{at the level of forms on the free word span first}, with operators descending only afterwards (an earlier draft quotiented first and used the contraction bound before proving it; that order is circular and is repaired here).  On the free vector space spanned by words define the sesquilinear form $\langle u, v\rangle := p(\bar u \circ v)$ (PSD by bond-RP; $\langle\varnothing,\varnothing\rangle = p(\varnothing) = 1$, so the vacuum is normalized) and the \emph{prepend} maps $S_x[v] = [x v]$.  \textbf{The letter must be prepended, not appended}: appending gives $\langle u, [v x]\rangle = p(\bar u\, v\, x)$, with the letter at the far end where site-RP says nothing — a false display in an earlier draft, twice caught in review and corrected here.  Prepending gives $\langle u, S_x v\rangle = p(\bar u\, x\, v)$, which by site-RP is a PSD form in $(u, v)$, and symmetric by reversibility ($p(\bar u x v) = p(\bar v x u)$).  Form inequalities, all on the free span: (i) $0 \le \langle v, S_x v\rangle$ and $\sum_x \langle v, S_x v\rangle = \langle v, T v\rangle$ with $T := S_0 + S_1$; (ii) for $\xi = \sum_i c_i [u_i]$, the quantity $F(k) := \langle\xi, T^k\xi\rangle = \sum_{i,j}\bar c_i c_j \sum_{|w| = k} p(\bar u_i\, w\, u_j)$ satisfies the uniform bound $F(k) \le \sum_{i,j}|c_i||c_j|$ (sums of word probabilities are $\le 1$) and the reflection Cauchy–Schwarz $F(k)^2 \le F(0)\,F(2k)$; iterating, $F(1) \le F(0)^{1 - 2^{-k}} F(2^k)^{2^{-k}} \to F(0)$, so $0 \le \langle\xi, T\xi\rangle \le \langle\xi, \xi\rangle$ as forms, and $0 \le \langle\xi, S_x\xi \rangle \le \langle\xi, T\xi\rangle$ likewise.  These bounds imply $S_x$ and $T$ annihilate null vectors of the form, hence descend to well-defined bounded self-adjoint operators $A_x, T$ on the quotient completion $\mathcal H$ with $0 \preceq A_x \preceq T \preceq \mathbf 1$ (in infinite dimensions, boundedness of the quotient operators is exactly what the form bounds provide, in the manner of Klein–Landau [6]).  Iterating the display, $\langle\omega, A_{x_1}\cdots A_{x_n}\omega\rangle = p(x_1\cdots x_n)$ with $\omega$ the class of the empty word; $T\omega = \omega$ is stationarity ($\sum_x p(xw) = p(w)$).  Finite Hankel rank bounds $\dim\mathcal H$ by the rank of the bond Gram. $\square$

\emph{Remark (what each reflection buys).}  Bond-RP alone builds the space and the Gram; \textbf{site-RP is exactly what makes the letters positive operators} — the two hypotheses are not interchangeable, and the conjunction is the honest name of the class this paper studies.  (An earlier draft derived $A_x \succeq 0$ from bond-RP "at a shifted bond"; that argument conflates the two reflections and is withdrawn.)

Thus RP processes are exactly the POVM-\emph{type} representable stationary processes (the letters are PSD but sum to a contraction $T \preceq \mathbf 1$, not to $\mathbf 1$ — the qualifier is kept deliberately): the realization conjecture asks when quantum-representable processes are secretly classical — a quantum-vs-classical boundary for stationary processes, related to but distinct from the hidden quantum Markov model literature [7] and to finitely correlated states [9], whose generating structure is the operator form above with the positivity demanded of the \emph{map} rather than the letters.

\textbf{Proposition 2.2 (exhaustion).}  For any data as in Theorem 2.1 — \emph{with or without} word positivity — every reflected Gram is PSD.

\emph{Proof.}  The display above never used positivity of $p$.  $\square$

Receipt: an engineered datum with $\min_w p(w) = -4.6\times 10^{-3}$ (word scan to length 10) has reflected Grams PSD at $-1.8\times 10^{-17}$.  Consequence: \emph{reflection is spent}.  All remaining force of the conjecture is the non-automatic requirement $\langle\omega, A_{x_1}\cdots A_{x_n}\omega\rangle \ge 0$ for non-commuting PSD factors.  We also document a trap this exposes: an apparent one-line proof of block-moment positivity via mismatched Gram families produces only the off-diagonal block of an automatically-PSD matrix, hence nothing.

\textbf{A genericity observation.}  For the natural construction (random PSD letters, $\omega$ the Perron vector of $T$), word positivity held in $160/161$ random draws across two studies; invalid instances had to be engineered by optimization.  Why the Perron state aligns so strongly with validity is open and flagged.

\section*{3. The clock, and that it fails RP}
\textbf{Example 3.1 (the clock, fully specified).}  A renewal-type quasi-realization on $\mathbb R^3$ [8]:

\[
\tau_1 = \tfrac12\begin{pmatrix} R(1) & 0\\ 0 & 1\end{pmatrix},
\qquad
\tau_0 = q\,r^{\mathsf T},\quad
q = \mathbf 1 - \tau_1\mathbf 1 =
\begin{pmatrix}1.150584\\ 0.309113\\ 0.5\end{pmatrix},\quad
r = \begin{pmatrix}0.07\\ 0\\ 0.93\end{pmatrix},
\]
with $R(1)$ rotation by $1$ radian; $p(w) = \pi\,\tau_{x_1}\cdots \tau_{x_n}\mathbf 1$ with $\pi = (0.038079, -0.021951, 0.983872)$ the stationary functional ($\pi\tau = \pi$, $\tau\mathbf 1 = \mathbf 1$, $\tau = \tau_0 + \tau_1$).  \textbf{Validity is proved to all orders in closed form}: because $\tau_0 = q r^{\mathsf T}$ has rank one, every word value factorizes over the runs of $1$s,

\[
p(1^{a_0}0\,1^{a_1}0\cdots 0\,1^{a_k})
= \bigl(\pi\tau_1^{a_0}q\bigr)
  \prod_{i=1}^{k-1}\bigl(r^{\mathsf T}\tau_1^{a_i}q\bigr)
  \cdot\bigl(r^{\mathsf T}\tau_1^{a_k}\mathbf 1\bigr),
\]
so validity is equivalent to nonnegativity of four scalar families of the form $\kappa^a(A + B\cos(a + \varphi))$ — the display shows single-$0$ separators, but the factorization covers \emph{all} words: run exponents $a_i = 0$ are allowed (consecutive zeros), the corresponding factor being $r^{\mathsf T} q = 0.5455 > 0$ — each family certified by its closed-form margin $A - |B| > 0$:

\begin{small}\begin{verbatim}
   family            A          |B|        margin A - |B|
   r  tau1^a q     0.465000   0.083397      0.381603
   pi tau1^a q     0.491936   0.052365      0.439571
   r  tau1^a 1     0.930000   0.098995      0.831005
   pi tau1^a 1     0.983872   0.062159      0.921714
\end{verbatim}\end{small}
The process is reversible, has Hankel rank exactly $3$, and its diagonal is

\[
p(1^n) = (1/2)^{n}\bigl(0.98387 + 0.06216\cos(n + \varphi)\bigr),
\]
an oscillation at one radian per step — irrational in turns — on the decay circle: by the obstruction theory for nonnegatively realizable sequences (the Berstel–Soittola theory; we cite the $\mathbb R_+$-rational version our application needs [10], not only the $\mathbb N$-rational statement: dominant singularities are the radius times roots of unity), \textbf{no finite positive realization exists at any dimension}. Functions of finite Markov chains are therefore a \emph{strict} subclass of valid reversible finite-rank processes.

\textbf{Proposition 3.2.}  The clock is not RP: the $15$-word reflected Gram over all words of length $\le 3$ has minimum eigenvalue $-1.06\times 10^{-2}$ (entries are finite products of the displayed rational-and-$R(1)$ data, so this is checkable to any precision; the negativity exceeds any conceivable rounding by ten orders), and the offset diagonal Hankel fails the moment test that Theorem 4.1 imposes on RP processes — consistent with the oscillating diagonal, which no moment sequence can produce.

This is the conjecture's first nontrivial confirmation: the unique known obstruction mechanism is incompatible with RP at the single-letter level.  The remainder of the paper is about \emph{block} diagonals, where the question is genuinely open.

\section*{4. Confinement: clocks need chiral blocks}
The companion paper proves: a word whose cyclic class is achiral factors (through a rotation) into two palindromes, hence its block operator is a product of two PSD matrices with real nonnegative spectrum, at every dimension.  Applied to processes:

\textbf{Theorem 4.1.}  Let $p$ be valid with the RP form (real form: real symmetric letters, real $\omega$ — available whenever $p$ is real, by realification), $W$ a block whose necklace is achiral. Then: (a) \emph{if the necklace contains a palindromic rotation} $W'$ (always true for odd $|W|$ — the reflection axis passes through a vertex — and for even $|W|$ exactly when some reflection axis is \textbf{edge–edge}, i.e. passes through two edge midpoints; e.g. $0110$ qualifies (edge–edge axis only), while $01$ and $0001$ do not (vertex–vertex axes only)).  An earlier draft stated the even-length criterion as "vertex–vertex" — that is exactly inverted, as the draft's own examples already betrayed, and we verified the corrected criterion exhaustively over all even achiral necklaces of length $\le 8$ ($45/45$, fixed-seed receipt).  At such a phase the diagonal $h(n) = p(W'^n)$ is an exact Stieltjes moment sequence ($B_{W'} \succeq 0$ and the spectral theorem); (b) at \textbf{every} phase of \textbf{every} achiral necklace — including those with no palindromic rotation, like $01$ — and for $n$ above the nilpotency index at $0$: $h(n) = \sum_k c_k\lambda_k^{\,n-1}$ with $\lambda_k \ge 0$ and real $c_k$ (exponent convention: $n - 1$ because $h(n) = \langle\omega', B_W^{\,n-1} B_W\,\omega\rangle$ groups one block with the boundary vectors; purely a normalization, stated to prevent an off-by-one reading) — the two-palindrome factorization makes $B_W$ a product of two PSD operators: real nonnegative spectrum, semisimple off zero — no oscillatory component; (c) $h$ admits a finite positive realization.  The hypotheses of the dominant-pole sufficiency of positive-realization theory [3, 4] are matched explicitly: (i) \emph{unique simple dominant pole} — the largest $\lambda_k$ with nonzero residue; if two $\lambda_k$ tie they merge into one residue, and ties cannot produce oscillation since all poles are nonnegative real; (ii) \emph{positive dominant residue} — forced by validity ($h(n) \ge 0$ with $h(n)/\lambda_{\max}^{\,n-1} \to c_{\max}$); (iii) \emph{degenerate cases}: if $h$ is eventually zero it is finitely supported and realized by a delay line outright; if every coupled residue vanishes, the effective dominant pole is the largest $\lambda$ surviving in $h$, and (i)–(ii) apply to it.  (iv) \emph{Absorbability lemma for the transient part}: the nilpotent component contributes finitely many initial terms $h(1), \ldots, h(\nu)$ — nonnegative by validity.  Let $(A, b, c)$ be the entrywise-nonnegative realization of the tail supplied by (i)–(iii), $h(\nu + k) = c^{\mathsf T} A^{k-1} b$ for $k \ge 1$.  The assembled system, displayed (on $\mathbb R^\nu \oplus \mathbb R^m$, with $S$ the $\nu$-stage shift $S_{k+1,k} = 1$ and $e_k$ coordinate vectors):

\[
x_{n+1} = \begin{pmatrix} S & 0 \\ b\,e_\nu^{\mathsf T} & A
\end{pmatrix} x_n, \qquad
y_n = \bigl(h(1), \ldots, h(\nu),\, c\bigr)^{\mathsf T} x_n,
\qquad x_1 = (e_1, 0),
\]
realizes $h$ exactly: for $n \le \nu$ the state is $(e_n, 0)$ and $y_n = h(n)$; at step $\nu$ the delay exit injects $b$ through the rank-one corner block, so $y_{\nu+k} = c^{\mathsf T} A^{k-1} b = h(\nu + k)$.  Every entry of the assembled matrices and the initial state is nonnegative, with no rescaling step (an earlier draft invoked an undefined "first-passage normalization" here; this displayed block system replaces it and is the justification of "absorbable").

\textbf{Corollary 4.2 (confinement).}  Any clock obstruction in a valid RP process lives on a chiral-necklace block; the smallest chiral classes have length six.

Receipts: Hankel matrices of palindromic-phase diagonals PSD at $-1.3\times 10^{-16}$ across 40 verified-valid processes; rotated phases show real poles ($|\mathrm{Im}| = 0$) with genuinely signed residues (to $-0.5$).

\section*{5. The chiral core: reduction and the normal escape}
\subsection*{5.1 The minimal dangerous configuration and its exact reduction}
A coupled complex pair strictly \emph{above} the coupled real spectrum violates diagonal positivity outright (equidistribution makes the cosine dominate), so the minimal threat is the \textbf{coupled triple}: top coupled spectrum $\{\rho, \rho e^{\pm i\theta}\}$, $\theta/2\pi$ irrational, semisimple.  (The spectral data $G, H, \beta, \alpha$ used throughout are defined once, in Theorem 5.1.)

\textbf{Theorem 5.1 (peripheral reduction, with the formulas).}  Write the triple's spectral data with right/left eigenvectors $R_s, L_s$ ($L^\dagger R = \mathbf 1$, $s \in \{0, +, -\}$), and set $G = L^\dagger L$, $H = G^{-1} = R^\dagger R$, $\beta_s = L_s^\dagger\omega$, $\alpha = G^{-1}\beta$, $D(x) = \mathrm{diag}(1, e^{ix}, e^{-ix})$.  Then for any alternating word $W^{a_1}\widetilde W^{b_1}W^{a_2}\cdots$ (at rank 3 exactly; at higher rank in the peripheral limit along exponent subsequences, with subperipheral corrections $o(\rho^{\Sigma})$):

\[
\rho^{-\Sigma}\,p\bigl(W^{a_1}\widetilde W^{b_1}\cdots\bigr)
= \bigl(D(a_1\theta)\,\bar\alpha\bigr)^{\mathsf T} G\,
  D(-b_1\theta)\,H\,D(a_2\theta)\,G\cdots(\text{final }\alpha
  \text{ or }\beta\text{ by parity}).
\]
Validity at all exponents plus Weyl equidistribution of $(a_j\theta, b_j\theta)$ forces these trigonometric forms to be nonnegative on the whole torus — a $\theta$-free system of conditions on the reduced datum $(G, \beta)$ alone (six real parameters in the gauge $G_{00} = G_{++} = G_{--} = 1$, $G_{0+}$ real).  \emph{Proof of the display}: insert the spectral resolution $\mathbf 1 = \sum_s R_s L_s^\dagger + (\text{subperipheral})$ between consecutive segments; each block power contributes $D(a\theta)$ on the peripheral part, each segment boundary contributes $G$ or $H$ by $L^\dagger R$-biorthogonality, and the subperipheral terms are $o(\rho^\Sigma)$ along the stated subsequences with \emph{fixed segment count} — the equidistribution argument is applied per fixed chain shape, with uniformity of the correction in the exponents of that shape (the scope sentence an earlier draft omitted).  The numerical validation at $1.1\times 10^{-11}$ over three- and four-segment words checks the implementation, not the proof.  Define for later use the \textbf{coupling ratio} $r := 2|c_+|/c_0$ with $c_s = \bar\alpha_s\beta_s$ (diagonal positivity gives $r \le 1$), and the \textbf{Gram distance} from normality $\mathrm{dist}(G) := \max(|G_{0+}|, |G_{+-}|)$ in the stated gauge — the two coordinates in which all Section 6 evidence is expressed.  The \textbf{subordination} $s$ of a (not necessarily peripheral) coupled complex pair is $|\lambda_{\mathrm{cx}}|/\rho_{\mathrm{coupled}}$, the pair's modulus relative to the coupled spectral radius.

Two structural discoveries inside this system.  First, \emph{entrywise} versus matrix positivity: the two-segment values $p(W^a\widetilde W^b)$ are probabilities, and their nonnegativity is a genuine constraint even though the corresponding Gram matrices are automatically PSD; probing it yields closed-form necessary inequalities (e.g. $\alpha_0^2 \ge 2|\alpha_+|^2(1 + |G_{+-}|)$ in a normalized gauge).  Second:

\textbf{Theorem 5.2 (normal escape).}  At $G = I$ — equivalently, at rank three, $B_W$ normal (at higher rank $G = I$ is normality \emph{of the peripheral part}; see the Section 6 scope correction) — every chain in the tower collapses to a single-frequency Fejér form: each chain value equals $c_0 + 2\,\mathrm{Re}(c_+ e^{i\Phi})$ with $\Phi$ the signed sum of segment phases.

\emph{Proof.}  $G = H = I$ makes all chain matrices diagonal; the boundary contractions then collapse the product to the displayed form.  $\square$

\textbf{Corollary 5.3.}  The all-depth torus system is satisfiable at every coupling ratio $r < 1$: take $G = I$ and any $\beta$ with $c_0 = 2|c_+|/r$ — every chain value is then $c_0(1 + r\cos\Phi) \ge c_0(1 - r) > 0$.  (Explicit witness used in the receipts: $G = I$, $\beta = (\sqrt2(1+10^{-6}), 1, 1)$, giving $r = 0.999998$ with all implemented depths' margins $+2\times 10^{-6}$ in the $c_0$-normalized convention — chain values divided by $c_0$; unnormalized, the floor is $c_0(1-r) \approx 4\times 10^{-6}$.)

\textbf{Consequence (scoped).}  The torus-constraint family of Theorem 5.1 — the complete set of conditions obtainable from words in $\{W, \widetilde W\}$ alone — does not separate normal-clock data from realizable data: any argument excluding the clock must use words leaving the block alphabet, or the letter structure itself. The question is thus pushed to (i) the \emph{letters} (can a chiral word of PSD letters be normal with nonreal spectrum? — the companion paper's Conjecture NR says no, with strong receipts), and (ii) \emph{depth dynamics off the normal point}.

\subsection*{5.2 Off the normal point: the deep adversary}
Each chain value is degree-one trigonometric in each phase separately, so exact cyclic coordinate descent gives a deep adversary (calibrated to reproduce the Fejér bound at $4\times 10^{-16}$ to depth 16).  Findings, reported as measured: a tower-optimized datum surviving grid constraints to depth 5 dies at depth 8 (margins $+0.018 \to -0.055 \to -0.338$ by depth 16): shallow feasibility is an artifact.  But band searches at fixed Gram distance $\varepsilon$ from $I$ give a mixed picture: some winners die at finite depth (instances at $k^* = 16$ and $48$), at least one survives a 64-phase adversary, and the $\varepsilon = 0.05$ shell retains coupling $0.84$ through constraint depth 56.  Whether all-depth feasibility is isolated at the normal point is open; stated as a conjecture in its own right:

\begin{quote}
\textbf{Conjecture ISO.}  If a reduced datum $(G, \beta)$ with nonzero coupling ($r > 0$) satisfies the all-depth torus system of Theorem 5.1, then $G = I$.
\end{quote}
Sum-of-squares certification over the torus is the named decisive tool.

\section*{6. The two-conjecture reduction and the evidence}
\textbf{Theorem 6.1 (conditional; rank-3 blocks).}  Assume (NR): a word in two PSD letters that is normal has real spectrum; and (ISO): all-depth torus feasibility with nonzero coupling forces $G = I$.  Then no valid RP process whose block operator $B_W$ has \textbf{rank three} carries a coupled irrational triple with semisimple peripheral part on that chiral block.

\emph{Proof.}  Such a process' reduced datum is all-depth feasible (Theorem 5.1), hence $G = I$ by (ISO); at rank three the peripheral triple \emph{is} the whole operator, so $G = I$ is normality of $B_W$, hence real spectrum by (NR) — contradicting the triple.  $\square$

\textbf{Scope correction (the higher-rank gap, found in review and kept visible).}  An earlier draft claimed the conclusion at every finite rank.  That does not follow: at rank $> 3$, $G = I$ makes the three \emph{peripheral} eigenvectors orthonormal among themselves but says nothing about their angle to the subperipheral spectral subspaces, so $B_W$ need not be normal and (NR) does not apply. Closing the general case needs one of: (ISO$^+$) all-depth feasibility forces peripheral-to-subperipheral orthogonality as well; or (NR$^{\mathrm{per}}$) a \emph{peripherally normal} variant of (NR) — a word operator whose peripheral part is normal has real peripheral spectrum.  Both are stated as named open strengthenings; the evidence below (walls, pinch) probes the reduced datum and is agnostic to the distinction, but the theorem itself is rank-3.

\textbf{Evidence, summarized.}  (a) \emph{Generic validity decouples}: across 120 verified-valid random processes, coupled complex poles never appear on any chiral block (subordination $s = 0$ throughout). (b) \emph{The wall}: steering searches under hard validity reach coupled complex poles but stall at $s \approx 0.19$–$0.21$ in three independent campaigns; every candidate at $s = 1$ fails validity at an explicit witness word — including a documented impostor whose first negative probability hides at word length $\approx 150$, beyond any letter scan (block-diagonal depth checks are therefore structural in all searches).  (c) \emph{The mechanism}: lifting a valid instance's complex pair to the coupled circle with residues fixed forces $h(n) < 0$ at $n = 32$.  (d) \emph{The pinch}: in the shared coordinate where $G = I$ is exactly normality, letter-realizable near-triples (modulus spread to $0.0015$, so the triple is kinematically reachable) live at Gram distance $0.50$–$0.94$ from normality, while deep-torus-feasible data live at $0.054$–$0.147$: \textbf{disjoint by a factor $\ge 3.4$} — a counterexample must thread two walls that no search has seen overlap.

\section*{7. Discussion}
\textbf{The assembly half.}  Even granting the spectral exclusion, realizing all block diagonals \emph{jointly} — one nonnegative pair on one cone — is the classical functions-of-Markov-chains problem [1, 2] and remains open; per-diagonal dominant-pole realizations construct different cones per block.  The natural program: assemble on the achiral sublattice first (exact Stieltjes structure provides canonical cones), then extend by reversal-pairing.

\textbf{Relation to quantum models.}  Theorem 2.1 places the conjecture at the boundary between POVM-representable and classically realizable processes.  Quantum trajectory models with non-normal Kraus operators \emph{can} oscillate at block level with automatic positivity, but generically fail bond-RP; the conjecture asserts that time-reflection symmetry is precisely what closes the quantum loophole at finite memory.

\textbf{Scope.}  Single irrational orbit plus real top, semisimple peripheral part; rational phases are unobstructed (block-power decimation); harmonically-related orbits need higher-degree Fejér systems (same machinery, unexecuted); peripheral Jordan tails open.

\section*{Reproducibility}
All quantitative claims regenerate from fixed-seed scripts: the representation and exhaustion receipts; the clock's certificates; the confinement receipts (40 valid processes); the reduction validation; tower searches with stated budgets; the adversary, band, and pinch measurements.  Code and canonical outputs accompany the submission; we flag explicitly that the Section 5.2/6 evidence (walls, bands, pinch, the 160/161 genericity study) is artifact-dependent — without the artifact those sections are claims, not evidence — so the artifact is part of the submission, not supplementary.

\section*{References}
[1] S. W. Dharmadhikari, Functions of finite Markov chains, \emph{Ann. Math. Statist.} \textbf{34} (1963) 1022–1032.

[2] A. Heller, On stochastic processes derived from Markov chains, \emph{Ann. Math. Statist.} \textbf{36} (1965) 1286–1291.

[3] B. D. O. Anderson, M. Deistler, L. Farina, L. Benvenuti, Nonnegative realization of a linear system with nonnegative impulse response, \emph{IEEE Trans. Circuits Syst. I} \textbf{43} (1996) 134–142.

[4] L. Benvenuti, L. Farina, A tutorial on the positive realization problem, \emph{IEEE Trans. Autom. Control} \textbf{49} (2004) 651–664.

[5] J. Glimm, A. Jaffe, \emph{Quantum Physics: A Functional Integral Point of View}, Springer (1987) — reflection positivity.

[6] A. Klein, L. J. Landau, Construction of a unique self-adjoint generator for a symmetric local semigroup, \emph{J. Funct. Anal.} \textbf{44} (1981) 121–137.

[7] A. Monras, A. Beige, K. Wiesner, Hidden quantum Markov models and non-adaptive read-out of many-body states, \emph{Appl. Math. Comput. Sci.} \textbf{3} (2011) 93.

[8] H. Jaeger, Observable operator models for discrete stochastic time series, \emph{Neural Computation} \textbf{12} (2000) 1371–1398.

[9] M. Fannes, B. Nachtergaele, R. F. Werner, Finitely correlated states on quantum spin chains, \emph{Commun. Math. Phys.} \textbf{144} (1992) 443–490.

[10] J. Berstel, C. Reutenauer, \emph{Noncommutative Rational Series with Applications}, CUP (2011), Ch. 8 — the $\mathbb R_+$-rational (Soittola) dominant-singularity theory used in Section 3.

\textbf{Companion paper:} \emph{Spectral reality of words in two positive semidefinite letters is a necklace invariant} (same author).

\end{document}
